\section{Actual factor descent and moving coefficient spaces}
\label{sec:coefficient-descent-v147}

We now separate the linear-preserver input, the descent of vector
bundles, and the recovery of coefficients.  The resulting principle
allows several varying coefficient spaces, of arbitrary prescribed
ranks, rather than only the line associated with a pencil.

\begin{lemma}[Linear preservers of the determinant cone]
\label{lem:linear-preserver-v147}
Let \(U,V,U',V'\) be complex vector spaces of dimension \(n\ge2\).
An invertible linear map
\(\Phi:\Hom(U,V)\to\Hom(U',V')\) carrying the determinant
hypersurface onto the determinant hypersurface has one of the forms
\[
 T\longmapsto gTh^{-1},\qquad
 T\longmapsto gT^{\mathsf t}h^{-1},
\]
where the second formula is written after choosing identifications
for transposition.  The two cases respectively preserve and exchange
the two rank-one rulings.  In the first case the pair \((g,h)\)
is unique up to a common nonzero scalar.
\end{lemma}
\begin{proof}
At a matrix of rank \(r\), an invertible \(r\)-minor reduces the
local rank condition to the rank condition on a generic
\((n-r)\)-square matrix.  The singular locus of the rank-at-most
\(s\) variety is therefore the rank-at-most \(s-1\) variety for
\(1\le s<n\): on the rank-\(s\) stratum the Schur-complement
entries give independent equations, whereas on a lower stratum the
linear terms of its defining \((s+1)\)-minors vanish.  Starting
with the determinant cone and iterating reduced singular loci hence
recovers the rank-one cone.  For \(n=2\) that cone is already the
determinant cone.

If \(v\otimes u\) and \(w\otimes z\) have rank one, their sum
has rank at most one only if \(v,w\) or \(u,z\) are dependent.
It follows that each maximal linear space of rank-one tensors has
one factor fixed.  The two families are
\(\Pj(V)\) and \(\Pj(U^*)\).  The induced Segre automorphism
must preserve these families or exchange them.  In the preserving
case its restriction to a ruling gives projective linear maps on
each factor.  Remove linear lifts of both maps.  The remaining
linear transformation fixes each decomposable line.  In tensor bases
its eigenvalues at \(v_i\otimes u_j\) agree along each row and
each column by applying it to
\((v_i+v_k)\otimes u_j\) and \(v_i\otimes(u_j+u_l)\).
They are all equal, so the remaining map is scalar.  This proves
the first form and its uniqueness.  Composing with transposition
reduces the exchanging case to the first.

This is the square invertible case of the classical rank-one
preserver theorem of Marcus and Moyls
\cite[Theorem~1, pp.~1218--1219]{MarcusMoyls}; the argument above
specifies the reduction from the determinant hypersurface that is
needed here.
\end{proof}

\begin{lemma}[An actual tensor map removes the line twist]
\label{lem:actual-factor-descent-v147}
Let \(B\) be a reduced complex variety, \(V\) a complex vector
space with \(\dim V\ge2\), and \(\mathcal A_1,\mathcal A_2\)
vector bundles.  Suppose an actual bundle isomorphism
\[
 \Psi:V\otimes\mathcal A_1\xrightarrow{\sim}V\otimes\mathcal A_2
\]
preserves the rank-one cones and induces the identity on the
\(\Pj(V)\)-parameter space of rulings.  There is a unique bundle
isomorphism \(a:\mathcal A_1\to\mathcal A_2\) with
\(\Psi=1_V\otimes a\).

Equivalently, the intrinsic operation of taking the commutant of the
constant matrix algebra gives
\begin{equation}\label{eq:commutant-descent-v147}
 \mathcal Hom_{\End(V)\otimes\OO_B}
       (V\otimes\mathcal A_1,V\otimes\mathcal A_2)
             \simeq\mathcal Hom_{\OO_B}(\mathcal A_1,\mathcal A_2).
\end{equation}
Thus after a constant left projective map has been lifted linearly,
no residual line-bundle twist is left in the recovery of the right
factor from \(\Psi\).
\end{lemma}
\begin{proof}
For every constant line \(\C v\subset V\), the corresponding
ruling is fixed, so on every geometric fibre \(\Psi\) maps
\(v\otimes\mathcal A_1\) onto \(v\otimes\mathcal A_2\).
Choose a basis \(v_1,\ldots,v_n\) for the purpose of checking
this assertion as an equality of bundle maps.  The rulings for
\(v_i\) force all off-diagonal blocks of \(\Psi\) to vanish.
The rulings for \(v_i+v_j\) force its diagonal blocks to agree.
Since \(B\) is reduced, fibrewise zero morphisms of vector bundles
are zero.  Therefore \(\Psi\) commutes with every constant
matrix unit, hence with \(\End(V)\otimes\OO_B\).

For clarity, commutant recovery itself is canonical and does not
use these chosen matrix units as extra data.  Tensoring a bundle
map \(a\) with \(1_V\) gives the forward map in
\eqref{eq:commutant-descent-v147}.  Commutation with matrix units
shows locally that it is bijective, with inverse reading any diagonal
block.  This proves that it is an isomorphism of sheaves, independent
of the checking basis.  Its inverse applied to the global section
\(\Psi\) produces a global section \(a\), not projective local
sections requiring arbitrary scalar lifts.  Applying the same
operation to \(\Psi^{-1}\) proves that \(a\) is invertible.
Uniqueness follows from injectivity of \(a\mapsto1_V\otimes a\).

Projective-factor data alone would permit replacing a bundle by a
line twist.  The additional datum used here is the actual linear map
of the tensor bundles.  It is this datum, recovered from the dual of
the first nilradical quotient, that makes the commutant argument
applicable.
\end{proof}

\subsection{A general finite-order coefficient inverse}
Fix \(n\ge2\), \(q\ge1\), and integers
\(0\le r_\lambda\le m_\lambda=\dim\mathbb S_\lambda\C^n\)
for partitions \(\lambda\) of \(q\) of length at most \(n\).
Assume that for at least one partition
\begin{equation}\label{eq:strict-support-v147}
                         0<r_\lambda<m_\lambda.
\end{equation}
Let \(B\) be a smooth connected projective complex variety,
\(V\) a constant \(n\)-dimensional space, and \(\mathcal U\)
a rank-\(n\) bundle.  For each \(\lambda\), choose a subbundle
with locally free quotient
\[
 \mathcal L_\lambda\subset\mathbb S_\lambda V^*\otimes\OO_B,
                       \qquad\rank\mathcal L_\lambda=r_\lambda.
\]
Zero and full subbundles are permitted.  In the matrix bundle
\(\mathcal E=\Hom(\mathcal U,V\otimes\OO_B)\), the
multiplicity-one Cauchy decomposition supplies the subbundle
\begin{equation}\label{eq:general-moving-coeff-v147}
 \mathcal J_q=\bigoplus_\lambda
       \mathcal L_\lambda\otimes\mathbb S_\lambda\mathcal U
             \ \subset\ \Sym^q\mathcal E^*.
\end{equation}
Let \(\mathcal D_n\subset\Sym^n\mathcal E^*\) be the
determinant line, put \(D=n+q\), and define
\begin{equation}\label{eq:coefficient-thickening-v147}
 \mathcal K=\mathcal D_n\mathcal J_q,
 \quad \mathcal A=\Sym\mathcal E^*/\bigl((\mathcal K)
                                  +\mathfrak a^{D+1}\bigr),
 \quad Y=\Spec_B\mathcal A,
\end{equation}
where \(\mathfrak a\) is the positive homogeneous ideal.  The
notation \(D\) in this subsection is an integer, not a failure
locus.  The scheme \(Y\) is finite over \(B\), and need not be
zero-dimensional over \(\C\).

\begin{theorem}[Unmarked descent for moving coefficient spaces]
\label{thm:general-moving-coefficients-v147}
For fixed \(n,q,(r_\lambda)\) satisfying
\eqref{eq:strict-support-v147}, the abstract complex scheme \(Y\)
recovers its coefficient data as follows.  Two such schemes \(Y_1,Y_2\)
are isomorphic if and only if there are an isomorphism
\(\sigma:B_1\to B_2\), a constant linear isomorphism
\(g:V_1\to V_2\), and an actual bundle isomorphism
\(h:\mathcal U_1\to\sigma^*\mathcal U_2\) such that, for every
\(\lambda\),
\begin{equation}\label{eq:coefficient-equivalence-v147}
 \mathcal L_{1,\lambda}=
       (\mathbb S_\lambda g)^*\sigma^*\mathcal L_{2,\lambda}.
\end{equation}
Every unmarked isomorphism induces a normal-bundle map
\(T\mapsto gTh^{-1}\); for that map the pair \((g,h)\) is
unique up to a common nonzero constant scalar.  Conversely every
such triple lifts homogeneously to an isomorphism.  The scheme is
finite locally free over its reduced base of rank
\[
             \binom{n^2+D}{D}-\sum_\lambda r_\lambda m_\lambda.
\]
No nontriviality hypothesis on \(\Pj_B(\mathcal U^*)\) is required.
\end{theorem}
\begin{proof}
\emph{Recovering the intrinsic linear data.}
Multiplication by a determinant polynomial is an injective linear
map of polynomial spaces.  In local bundle frames its matrix is
constant over \(\C\) and its cokernel is free.  The Cauchy
summands and the \(\mathcal L_\lambda\) are subbundles, so
\(\mathcal K\) is a subbundle of rank
\(\sum r_\lambda m_\lambda\).  The graded pieces of
\(\mathcal A\) below \(D\) are the full symmetric powers;
its top piece is \(\Sym^D\mathcal E^*/\mathcal K\).
This proves local freeness and the rank formula.

The nilradical \(\mathfrak n\) of \(\OO_Y\) is its positive
ideal, since \(B\) is reduced.  Thus \(Y_{\mathrm{red}}=B\) and
\(\mathfrak n/\mathfrak n^2=\mathcal E^*\) intrinsically.
The multiplication kernel
\[
 \ker\bigl[\Sym^D(\mathfrak n/\mathfrak n^2)
                                    \longrightarrow\mathfrak n^D\bigr]
\]
is precisely \(\mathcal K\), by the proof of
Proposition~\ref{prop:first-relation-orbit-v147} in local frames.
An abstract isomorphism consequently gives the base isomorphism
\(\sigma\) and an actual linear normal-bundle isomorphism carrying
\(\mathcal K_1\) to \(\mathcal K_2\) contravariantly.  The
normal map is linear over the reduced base even when the original
scheme map does not preserve the chosen graph projection.  Indeed
the action on \(\mathfrak n/\mathfrak n^2\) is linear over
\(\OO_Y/\mathfrak n\).

\emph{Recovering and orienting the factors.}
Generate the homogeneous ideal \((\mathcal K)\) in the recovered
symmetric algebra; do not confuse its cone with the finite
truncation.  Its radical is the determinant ideal.  To verify this,
at an invertible matrix \(T\), the map
\(\mathbb S_\lambda T\) is invertible.  For a partition with
\(r_\lambda>0\), a nonzero covector in \(\mathcal L_\lambda\)
has a nonzero matrix coefficient there.  Thus the residual
coefficients cannot all vanish off the determinant.  All generators
vanish on the determinant.  The relative determinant hypersurface
is reduced over the smooth base; the equality of radicals follows
on affine charts from this equality of closed points.  Successive
reduced singular loci and Lemma~\ref{lem:rank-one-fano-scheme}
recover its two unordered tensor rulings with their relative scheme
structures.

Cancellation of the determinant line is intrinsic and recovers the
ideal generated by \(\mathcal J_q\), and hence its degree-\(q\)
piece.  In a Cauchy summand the two contraction supports of
\(\mathcal L_\lambda\otimes\mathbb S_\lambda\mathcal U\)
have ranks \((r_\lambda,m_\lambda)\).  Separate factor changes
preserve these ranks.  By Lemma~\ref{lem:linear-preserver-v147},
an exchange of the two matrix factors interchanges the ranks in
each summand.  At a partition satisfying
\eqref{eq:strict-support-v147}, the result
\((m_\lambda,r_\lambda)\) cannot equal
\((r_\lambda,m_\lambda)\) for the other datum.  Thus no fibre
permits an exchange.  Twisting local factor lifts by lines changes
no support rank.  The \(V\)-ruling is thereby oriented intrinsically,
including when both projective ruling bundles are trivial.

\emph{Descending the actual map.}
By Lemma~\ref{lem:oriented-segre-descent}, the induced morphism
\(B_1\to\operatorname{Isom}(\Pj(V_1),\Pj(V_2))\) is constant:
the target is affine and \(H^0(B_1,\OO_{B_1})=\C\).  Choose a
constant linear lift \(g\).  Remove it from the actual normal map.
Lemma~\ref{lem:actual-factor-descent-v147}, applied to
\(\mathcal A_1=\mathcal U_1^*\) and
\(\mathcal A_2=(\sigma^*\mathcal U_2)^*\), supplies the unique
global right factor.  Dualizing its inverse gives the actual
isomorphism \(h\).  This step uses more than an isomorphism of
projective bundles; it is exactly what excludes an undetected line
twist.  Replacing \(g\) by \(cg\) replaces \(h\) by \(ch\).
These are all ambiguities because the left lift is constant.

\emph{Recovering every coefficient subbundle.}
The oriented Cauchy decomposition is now functorial for the recovered
left-right map.  Project the residual degree-\(q\) module onto each
summand and contract against the full right dual factor.  The image
is precisely \(\mathcal L_\lambda\), as a subbundle of the
constant left module, not just a collection of fibrewise projective
subspaces.  The right map acts invertibly on the full Schur factor
and hence does not change this image.  Pullback by the normal map
therefore gives \eqref{eq:coefficient-equivalence-v147} for the same
constant \(g\) in all summands.  Scalar choices in the determinant
line do not change these subbundles.  This proves necessity.

Conversely \(T\mapsto gTh^{-1}\) carries the determinant line,
all coefficient subbundles, and every power of the positive ideal
to the corresponding objects.  It induces the required homogeneous
scheme isomorphism.  The stated uniqueness refers to the induced
normal map; it does not eliminate the higher-order automorphisms of
that scheme.
\end{proof}

\subsection{Identical fibres and graded bundles do not determine a family}
The global conclusion has content even when neither fibre isomorphism
classes nor the underlying graded vector bundles vary.  The following
example uses the general coefficient theorem, not a further spectral
readout of a reconstructed pencil.

\begin{example}[Two isotrivial thickenings separated by coefficient descent]
\label{ex:isotrivial-covers-v147}
Take \(B=\Pj^1\), \(V=\C^3\), \(\mathcal U=\OO_B^3\),
\(q=1\), and \(r_{(1)}=1\).  The two everywhere nonzero columns
\[
 (s^3,t^3,0)^{\mathsf t},\qquad
 (s^3+st^2,t^3,0)^{\mathsf t}
\]
define subbundles \(\mathcal L_0,\mathcal L_1\subset V^*\otimes\OO_B\),
both isomorphic to \(\OO_B(-3)\).  Let \(Y_0,Y_1\) be their
order-\(4\) thickenings \eqref{eq:coefficient-thickening-v147}.
They are finite flat of rank \(712\) over \(B\) and Zariski
locally trivial with the same Artin algebra as fibre.  Their
nilradical graded vector bundles are isomorphic in every degree,
although \(Y_0\) and \(Y_1\) are not isomorphic as abstract
complex schemes.
\end{example}
\begin{proof}
Neither column vanishes, since its second entry is nonzero when
\(t\ne0\), and its first is nonzero at \(t=0\).  The left
coefficient map to its image line \(\Pj^1\subset\Pj(V^*)\)
is, in the coordinate \(z=s/t\), respectively
\[
                        f_0(z)=z^3,\qquad f_1(z)=z^3+z.
\]
A local frame completing a generator of \(\mathcal L_i\) to a
basis of \(V^*\otimes\OO\) changes the coefficient line to a
fixed line.  This gives Zariski local product presentations of both
thickenings.  Alternatively, equality of the fibre isomorphism
classes follows from transitivity of \(\operatorname{GL}(V)\)
on nonzero coefficient lines.  The rank formula is
\(\binom{13}{4}-3=712\).

Let \(E=V^*\otimes\C^3\), of dimension nine.  For \(j<4\),
the degree-\(j\) nilradical quotient is \(\Sym^j E\otimes\OO_B\).
Multiplication by the fixed determinant embeds \(E\) as a fixed
nine-dimensional subspace of \(\Sym^4 E\).  Choosing a constant
linear complement gives, for the top quotient,
\[
 \mathfrak n_i^4\simeq
 \OO_B^{486}\oplus
           \bigl((V^*\otimes\OO_B)/\mathcal L_i\bigr)^{\oplus3}.
\]
Each column has only its first two entries nonzero, and these two
entries have no common zero.  The rank-one quotient of \(\OO_B^2\)
by \(\OO_B(-3)\) is consequently \(\OO_B(3)\), by determinants.
Thus the displayed top bundle is
\(\OO_B^{489}\oplus\OO_B(3)^{\oplus3}\) for both \(i\).
The assertions about every graded vector bundle follow.  They are
not assertions that the graded \emph{algebra} structures agree.

If an abstract isomorphism \(Y_0\simeq Y_1\) existed,
Theorem~\ref{thm:general-moving-coefficients-v147} would identify
their coefficient maps by one automorphism of \(B\) and one
constant projective transformation of \(V^*\).  Both maps have
as image the same projective line, so the latter transformation
would restrict to a projective automorphism of that line.  Hence
\(f_0\) and \(f_1\) would be equivalent by independent projective
changes in source and target.  This is impossible.  The map \(f_0\)
has two ramification points, each of index three.  The map \(f_1\)
has the two simple zeros of \(3z^2+1\), each of index two, and
infinity, of index three.  Ramification indices are invariant under
those projective changes.  Thus the unmarked algebra scheme
distinguishes these two globally different coefficient maps even
though its fibre algebras, conormal bundle, abstract first-relation
bundle, and all graded vector bundles agree.
\end{proof}

\subsection{Pencils and the category of their automorphisms}
For pencils, take \(q=2p\), \(p=N-2\), and
\(F(H)=\bigwedge^{N-2}\Sym^2 H\).  By
Lemma~\ref{lem:pencil-irreducibility} this is a single Schur module
of dimension \(M=\binom N2\), and
Lemma~\ref{lem:moving-coefficients-v146} identifies its coefficient
subbundle with
\(\mathcal L_{\mathcal R}=\bigwedge^p(\Sym^2V/\mathcal R)^*\).
Its rank is one, so the strict-support hypothesis holds for every
pencil.  Exterior duality recovers the subbundle \(\mathcal R\)
from that line.  Theorem~\ref{thm:general-moving-coefficients-v147}
therefore proves the unrestricted moving-pencil inverse, including
its actual source-bundle conclusion.  This is a simultaneous
reconstruction theorem for a map into a coefficient Grassmannian,
not just a classification of its individual fibres.

\begin{corollary}[Algebraic automorphism sequence for a pencil]
\label{cor:pencil-group-scheme-v147}
Let \(\mathfrak A_R^{[d]}\) be the unmarked local algebra of
Theorem~\ref{thm:artin-local-inverse-v146}, and let
\(G_R=\{g\in\operatorname{GL}(V):(\Sym^2g)R=R\}\).
There is a split exact sequence of complex affine algebraic groups
\begin{equation}\label{eq:pencil-algebraic-group-v147}
 1\longrightarrow U_R\longrightarrow
 \operatorname{Aut}(\mathfrak A_R^{[d]})
 \longrightarrow (G_R\times\operatorname{GL}(U))/\mathbb G_m
 \longrightarrow1,
\end{equation}
where \(\mathbb G_m\) is the diagonal scalar subgroup.  Its
intrinsic quotient is the closed stabilizer of \(K_R\) in the
general linear group of the cotangent space.  In that form the
sequence represents automorphisms after extension to every
commutative complex algebra.  Its splitting is homogeneous after
choosing generators, and \(U_R\) has the intrinsic description and
dimension in Proposition~\ref{prop:first-relation-group-v147}.
\end{corollary}
\begin{proof}
Proposition~\ref{prop:first-relation-group-v147} gives the represented
sequence with quotient \(H_{K_R}\).  The coefficient inverse and
Lemma~\ref{lem:linear-preserver-v147} identify its complex points
with the image of the left-right action; transposition is excluded
by the support ranks \((1,M)\).  Use the contragredient of the
normal action to obtain a homomorphism on the cotangent space.
Its kernel is exactly the diagonal scalar subgroup: a tensor product
acting identically has both factors scalar, with reciprocal actions
on the dual factors.  The image is a closed algebraic subgroup and
all the groups in question are smooth in characteristic zero
\cite[Lemma~39.8.2]{StacksGroups}.  Thus the stabilizer and this
image agree as closed subgroup schemes, not only on complex points.
The induced quotient identification is an isomorphism of algebraic
groups: after division by the kernel it is a bijective separable
homomorphism between smooth groups.  The splitting and kernel are
those already constructed by substitutions.
\end{proof}

For a positive-dimensional projective reduction,
Theorem~\ref{thm:automorphism-kernel-v146} is used as an exact
sequence of groups of complex scheme automorphisms, with its stated
filtered kernel.  We do not identify that sequence with a finite-type
affine group scheme or a moduli-stack equivalence.  Likewise, the
marked construction over nonreduced parameter bases later in the
paper retains its specified relative socle: it does not replace that
ideal by the absolute nilradical.  These distinctions preserve the
full unmarked geometric assertion while specifying its category.
